\documentclass[10pt,a4paper]{amsart}
\usepackage[margin=19mm]{geometry}
\usepackage{amsmath,amssymb,mathtools}
\usepackage[scr=rsfs,cal=euler]{mathalfa}
\usepackage{xfrac}
\usepackage[unicode,hidelinks,bookmarks=false]{hyperref}
\newcommand{\ie}{i.e.\ }
\newcommand{\cf}{cf.\ }
\newcommand{\resp}{respectively\ }
\newcommand{\Subsection}{\subsection}
\providecommand{\nobreakdash}{\nobreak\hbox{-}}
\setlength{\emergencystretch}{2em}
\multlinegap0pt
\usepackage{natbib}
\newtheorem{remark}{Remark}
\newtheorem{theorem}{Theorem}
\title{Metric basis and dimension of barycentric subdivision of zero divisor graphs}
\author{S. Vidya, Sunny Kumar Sharma, Prasanna Poojary, Omaima Alshanqiti,, G. R. Vadiraja Bhatta}
\date{Source: arXiv:2602.11816v1 (CC BY 4.0)}
\begin{document}
\maketitle

\noindent\emph{Source and licence.} Metric basis and dimension of barycentric subdivision of zero divisor graphs. Version arXiv:2602.11816v1 (CC BY 4.0), distributed by arXiv under the Creative Commons Attribution 4.0 International licence, http://creativecommons.org/licenses/by/4.0/, as recorded in the arXiv licence metadata for this exact version and shown on the abstract page https://arxiv.org/abs/2602.11816v1. Full licence notice: \texttt{licenses/arxiv-2602.11816-CC-BY-4.0.txt}.\par\smallskip

\noindent\textit{Editorial scope.} Counted unit: the article's theorem theorem (source0.tex lines 361--364). The article's complete original proof of it is delivered with it (source0.tex lines 365--426, one \texttt{proof} environment of 62 lines). No mathematics is written, repaired or re-derived: the statement and the proof are the article's own lines, in the article's own notation, and no retained line is substituted or re-flowed.  Retained uncounted as the source-local context the proof needs: lines 198--203; lines 351--353; lines 354--359.  One declared adaptation: exactly 1 ancillary strings inside the retained spans are omitted (image-inclusion commands and, where the source repeats a label, the repeated label definitions) and no other line differs from the source; the figure environments, with their placement, captions and labels, are kept, and the package records the omitted strings in fidelity receipts bound to the affected bodies.  No retained line contains a citation, so the wrapper carries no bibliography and no bibliography entry.  The retained lines contain the article's own diagram code, which the wrapper typesets with the standard tikz packages; no image file is delivered and the package declares no asset dependency.  Editorial formatting only, in addition: an AMS \texttt{amsart} preamble that restates the article's own declarations and macros used by the retained lines, namely the environments \texttt{remark}, \texttt{theorem} on separate counters, together with the standard packages those lines need.  The retained environments print in this document as Remark 1, Theorem 1 in order of appearance, whereas the article prints the same items with the numbering of its own full document.  The complete native source of the article as distributed in the arXiv e-print for this exact version is bundled unchanged as \texttt{sources/194494/source0.tex}.
\begin{figure}[ht!]
    \centering
    
    \caption{Barycentric Subdivision of Zero Divisor Graph of $\mathbb{Z}_{pq}$}
    \label{fig:4}
\end{figure}

\begin{remark}\label{22}
  From Figure \ref{fig:4}, it is evident that, if $d(a_i,x)\neq 1$ and $d(a_j,x) \neq 1$, where $1\leq i, j \leq q-1$, $i\neq j$, and $x\in V(BS(\Gamma(\mathbb{Z}_{pq})))\backslash\{a_i,a_j\}$, then $d(a_i,x)=d(a_j,x)$.     
 \end{remark}

 \begin{remark} \label{last} \quad \\
    \begin{itemize}
        \item From Figure \ref{fig:4}, it is clear that among the \( p-1 \) vertices of the set \( Q \), exactly \( \frac{p-1}{2} \) vertices are adjacent only to the vertices of the set \( B \), and not to any vertices in other sets. Similarly, exactly \( \frac{p-1}{2} \) vertices are adjacent only to vertices of the set \( C \), and not to any vertices in other sets.
    \item From Figure \ref{fig:4}, it is evident that, if $d(q_i,a)\neq 1$ and $d(q_j,a) \neq 1$, where $1\leq i, j \leq p-1$, $i\neq j$, and $a\in V(BS(\Gamma(\mathbb{Z}_{pq})))\backslash\{q_i,q_j\}$, then $d(q_i,a)=d(q_j,a)$.\\     
\end{itemize}
\end{remark}

\begin{theorem}
Let $p$ and $q$ be two distinct odd primes, where $p\geq 5$ and $p+1<q<2p-1$. Suppose $R=\mathbb{Z}_n$, where $n=pq$, then
$dim(BS(\Gamma(R))) > q-2$. 
 \end{theorem}

    \begin{proof}
\noindent Consider the ring $R = \mathbb{Z}_{pq}$, where $p \geq 5$ and $p + 1 < q < 2p - 1$. Let $\Gamma(R)$ represent the zero divisor graph of this ring. We denote $BS(\Gamma(R))$ as the barycentric subdivision of $\Gamma(R)$. The graph $BS(\Gamma(R))$ consists of $pq-1$ vertices, and $2(p-1)(q-1)$ edges. In particular, we partitioned the vertex set of $BS(\Gamma(R))$ into 4 sets namely $A=\{a_1,a_2,\ldots, a_{q-1}\}$,  $B=\{b^\nu_i\}$, $C=\{c^\nu_i\}$, where $1\leq \nu\leq \frac{p-1}{2}$, and $1\leq i\leq {q-1}$ ($B^\nu=\{b^\nu_1,b^\nu_2,\ldots,b^\nu_{q-1}\}$, $C^\nu=\{c^\nu_1,c^\nu_2,\ldots,c^\nu_{q-1}\}$, where $1\leq \nu \leq \frac{p-1}{2}$ are the subsets of $B$ and $C$, respectively), and $Q=\{q_1,q_2, \ldots, q_{p-1}\}$. These sets are illustrated in Figure \ref{fig:4}. \\

\noindent Now we want to prove that dim($BS(\Gamma(R)) > q-2$ when $p+1< q<2p-1$. Let us assume that $W$ is a resolving set of $BS(\Gamma(R))$, and $|W|=q-2$. Then we get the following contradictions. \\


\begin{enumerate}
\item If $q-2$ vertices of the set $W$ are chosen from the set $A$, then, according to Remark \ref{last}, all the vertices of $Q$ have the same metric coordinates with respect to $W$, i.e., $\delta(q_{u}|W) =\delta(q_{v}|W)$, where $q_{u},q_{v} \in Q.$
\item If $q-2$ vertices of the set $W$ are chosen from $B$ (or $C$), then, according to Remark \ref{last}, there will exist at least two vertices of $Q$ that are not adjacent to any vertices of $W$, say $q_{u}$ and $q_{v}$, and those two have the same metric coordinates with respect to $W$ i.e., $\delta(q_{u}|W) =\delta(q_{v}|W)$.
% \item Choose $q-2$ vertices from  $C$, then
% by remark \ref{last}, we can deduce the existence of at least two vertices of $Q$ that are not adjacent to any vertices of $W$, say $q_{u}$, and $q_{v}$, and those two have the same metric coordinates, i.e., $\delta(q_{u}|W) =\delta(q_{v}|W)$.
\item If $q-2$ vertices of $W$ are chosen from  $A$ and $B$ (selecting at least one vertex from each set), then
from Remark \ref{last}, there will exist at least two vertices of $Q$ that are adjacent only to vertices of $C$, say $q_{u}$ and $q_{v}$, and those two have same metric coordinates, i.e., $\delta(q_{u}|W) =\delta(q_{v}|W)$.
\item If $q-2$ vertices of $W$ are chosen from $A$ and $C$, then
from Remark \ref{last}, there will exist at least two vertices of $Q$ that are adjacent only to vertices of $B$, say $q_{u}$ and $q_{v}$, and those two have same metric coordinates, i.e., $\delta(q_{u}|W) =\delta(q_{v}|W)$.
\item  \label{item1} If $q-2$ vertices of $W$ are chosen from \( A \) and \( Q \); then, there will exist at least two vertices in \( A \), denoted as \( a_u \) and \( a_v \), where \( 1 \leq u, v\leq q-1\) and \(u\neq v \), and these two vertices are not adjacent to any vertex of  \( W \) and do not belong to set \( W \). This is because from the fact that \( |A| = q-1 \), and the neighbourhoods of distinct vertices in \( A \) are disjoint, i.e., \( N[a_r] \cap N[a_s] = \emptyset \) for \( 1 \leq r, s \leq q-1 \) with \( r \neq s \), and also one vertex is choosen from \( Q \), and all vertices in \( A \) are equidistant from every vertex in \( Q \), it follows (as stated in Remark \ref{22}) that  \( a_u \) and \( a_v \) have the same metric coordinates with respect to \( W \), i.e., $\delta(a_u | W) = \delta(a_v | W)$.
\item If $q-2$ vertices of $W$ are chosen from  $B$ and $C$, then we have the following cases: 
\begin{itemize}
    \item If at least two sets, say \(B^{\nu_1}\) and \(B^{\nu_2}\) (or \(C^{\nu_1}\) and \(C^{\nu_2}\), or \(B^{\nu_1}\) and \(C^{\nu_3}\), where \(1 \leq \nu_1, \nu_2, \nu_3 \leq \frac{p-1}{2}\), and \(\nu_1 \neq \nu_2\)), have exactly one vertex of \(W\), while the remaining sets ($B^{\nu_3}$, $C^{\nu_4}$, where \(1\leq \nu_4, \nu_3 \leq \frac{p-1}{2}\), and $\nu_1 \neq \nu_3 \neq \nu_2$) have either more than one vertex or do not have any vertices of \(W\), then we encounter the following contradictions:
    
(without loss of generality, consider the sets \(B^{\nu_1}\) and \(B^{\nu_2}\) as having exactly one vertex of \(W\), say \(b^{\nu_1}_i\) and \(b^{\nu_2}_j\), respectively, where \(1 \leq i, j \leq q-1\).)
    \begin{enumerate}
        \item If each $a_k$ ($1\leq k \leq q-1$) adjacent to at most one vertex of $W$, then $\delta( b^{\nu_2}_i|W) =\delta( b^{\nu_1}_{j}|W)$.
\item If at least one vertex in $A$, say $a_k$ ($1\leq k \leq q-1$) is adjacent to more than one vertex of $W$, then there will exist at least two vertices $a_u$ and $a_v$ in $A$, where $1 \leq u \neq v \leq q-1$, that are not adjacent to any vertex of $W$, and do not belong to set $W$. This is because $|A|=q-1$, and $N[a_r]\cap N[a_s]=\emptyset$, where $1\leq r,s\leq q-1$, $r\neq s$. According to Remark \ref{22}, these two vertices $a_u$ and $a_v$ have the same metric coordinates with respect to $W$, i.e., $\delta(a_u|W) = \delta(a_v|W)$.
 \end{enumerate}
    \item If only one set, say $B^{\nu_1}$ (or $C^{\nu_2}$), where $1 \leq \nu_1, \nu_2 \leq \frac{p-1}{2}$, has one vertex of $W$, say $b^{\nu_1}_i$, while the remaining sets have either more than one vertex of $W$ or do not have any vertices of $W$, then we encounter the following contradictions:
    \begin{enumerate}
        \item If each $a_k$ ($1\leq k \leq q-1$)  adjacent to at most one vertex of $W$, then $\delta(a_i|W) = \delta(q_j|W)$, where $q_j$ and $a_i$ is adjacent to $b^{\nu_1}_i$.
        \item If at least one vertex in $A$, say $a_k$ ($1\leq k \leq q-1$) is adjacent to more than one vertex of $W$, then there must be at least two vertices $a_u$ and $a_v$ in set $A$ (where $1 \leq u \neq v \leq q-1$) that are not adjacent to any vertex of $W$ and do not belong to set $W$. This is because $|A|=q-1$, and $N[a_r]\cap N[a_s]=\emptyset$, where $1\leq r,s\leq q-1$, $r\neq s$. According to Remark \ref{22}, these two vertices $a_u$ and $a_v$ have the same metric coordinates with respect to $W$, i.e., $\delta(a_u|W) = \delta(a_v|W)$.
        \end{enumerate}
     \item If all the sets $B^{\nu_1}$ and $C^{\nu_2}$ (where $1\leq \nu_1,\nu_2 \leq \frac{p-1}{2})$ either contain more than one vertex of  $W$ or do not contain any vertex of $W$, then there exist at least two sets do not contain any vertex of $W$, which we will denote by $B^{\nu_3}$ and $B^{\nu_4}$. This is because $W$ has $q-2$ vertices, and $q-2<2p-4$. Therefore, From Remark \ref{last} $\delta(q_i|W) = \delta(q_j|W)$, where $q_i$ and $q_j$ are adjacent to vertices of $B^{\nu_3}$ and $B^{\nu_4}$ respectively.  
   \end{itemize}
\item If $q-2$ vertices of $W$ are chosen from $B$ and $Q$ (selecting at least one vertex from each set), then $\delta(a_i|W) = \delta(a_j|W)$, where $a_i$ and $a_j$ are not adjacent to any vertices of $W$. The reason is the same as in the case \ref{item1}.
 \item If $q-2$ vertices of $W$ are chosen from $C$ and $Q$ (selecting at least one vertex from each set), then $\delta(a_i|W) = \delta(a_j|W)$, where $a_i$ and $a_j$ are not adjacent to any vertices of $W$. The reason is the same as in the previous case \ref{item1}.
 \item If $q-2$ vertices of $W$ are chosen from $A$, $B$, and $C$, then we have the following cases
 \begin{itemize}
     \item  If we choose at least two vertices from $A$, say $a_{k_1}$ and $a_{k_2}$, then
     \begin{enumerate}
         \item If at least two sets, say $B^{\nu_1}$ and $B^{\nu_2}$ (or $C^{\nu_1}$ and $C^{\nu_2}$, or $B^{\nu_1}$ and $C^{\nu_3}$), where $1\leq \nu_1, \nu_2, \nu_3 \leq \frac{p-1}{2}, \nu_1\neq \nu_2$, have exactly one vertex of $W$, say $b^{\nu_1}_i$ and $b^{\nu_2}_j$ ($1\leq i,j \leq q-1)$, respectively and remaining sets have more than one or do not have any vertices of $W$, then we encounter the following contradiction: 
         \begin{enumerate}
        \item If each $a_k$ ($1\leq k \leq q-1$) is adjacent to at most one vertex of $W$ except $a_{k_1}$ and $a_{k_2}$, and $a_{k_1}$ and $a_{k_2}$ not adjacent to any vertex of $W$, then $\delta( b^{\nu_2}_i|W) =\delta( b^{\nu_1}_{j}|W)$.
       \item If at least one vertex in $A$, say $a_k$ ($1\leq k \leq q-1$) is adjacent to more than one vertex of $W$, or If at least one vertex in $A$ say $a_k$ ($1\leq k \leq q-1$) is belonging to set $W$, and adjacent to at least one vertex of $W$, then there will exist at least two vertices in $A$ , say $a_u $ and $a_v $, where $1\leq u \neq v \leq q-1$, not adjacent to any vertex of $W$, and also not belonging to set $W$. This is because $|A|=q-1$, and $N[a_r]\cap N[a_s]=\emptyset$, where $1\leq r,s\leq q-1$, $r\neq s$. From Remark \ref{22} those two vertices $a_u $ and $a_v $ have same metric coordinates with respect to $W$, i.e.,      
 $\delta(a_{u}|W) =\delta(a_{v}|W)$.
    \end{enumerate}
    \item \label{item 9} If only one set, say \( B^{\nu_1} \) (or \( C^{\nu_2} \)), where \( 1 \leq \nu_1, \nu_2 \leq \frac{p-1}{2} \), contains one vertex of  \( W \), say \( b^{\nu_1}_i \), while the remaining sets contain either more than one or do not have any vertices of  \( W \), then there exist at least two sets do not contain any vertex of $W$, which we will denote by $B^{\nu_3}$ and $B^{\nu_4}$ (where $\nu_3 \neq \nu_4$). This is because $W$ has $q-2$ vertices, and $q-2<2p-4$, and also 2 vertices of $W$ are chosen in $A$. Therefore, From Remark \ref{last} $\delta(q_i|W) = \delta(q_j|W)$, where $q_i$ and $q_j$ are adjacent to vertices of $B^{\nu_3}$ and $B^{\nu_4}$ respectively.
 \item If all sets $B^{\nu_i}$ and $C^{\nu_j}$ (where $1\leq i,j \leq \frac{p-1}{2})$ either contain more than one vertex of $W$ or do not contain any vertex of $W$, then we arrive at the same contradiction mentioned in the previous case \ref{item 9}. 
   \end{enumerate} 
    \item If we choose 1 vertex from $A$, say $a_{k_1}$ ($1\leq k_1 \leq q-1$), and remaining vertex from $B$ and $C$, then we have the following cases
    \begin{enumerate}
        \item If at least two sets, say $B^{\nu_1}$ and $B^{\nu_2}$ (or $C^{\nu_1}$ and $C^{\nu_2}$, or $B^{\nu_1}$ and $C^{\nu_3}$, where $1\leq \nu_1,\nu_2,\nu_3 \leq \frac{p-1}{2}, \nu_1\neq \nu_2)$), have exactly one vertex of $W$, say $b^{\nu_1}_i$ and $b^{\nu_2}_j$ (where $1\leq i,j \leq \frac{p-1}{2})$, respectively and remaining sets have more than one or do not have any vertices of $W$, then we get the following contradictions:
         \begin{enumerate}
        \item If each $a_k$ ($1\leq k \leq q-1$) is adjacent to at most one vertex of $W$ except $a_{k_1}$, and $a_{k_1}$ is not adjacent to any vertex of $W$, then $\delta( b^{\nu_2}_i|W) =\delta( b^{\nu_1}_{j}|W)$.
        \item If at least one vertex in $A$, say $a_k$ ($1\leq k \leq q-1$) is adjacent to more than one vertex of $W$, or if one vertex in $A$, say $a_{k_1}$ ($1\leq k_1 \leq q-1$) is belonging to set $W$ and adjacent to at least one vertex of $W$, then there must be at least two vertices $a_u$ and $a_v$ in set $A$, where $1 \leq u \neq v \leq q-1$, that are not adjacent to any vertex of $W$ and do not belong to set $W$. This is because $|A|=q-1$, and $N[a_r]\cap N[a_s]=\emptyset$, where $1\leq r,s\leq q-1$, $r\neq s$. According to Remark \ref{22}, these two vertices $a_u$ and $a_v$ have the same metric coordinates with respect to $W$, i.e., $\delta(a_u|W) = \delta(a_v|W)$.
        \end{enumerate}
    \item  If only one set say, $B^{\nu_1}$  (or $C^{\nu_2}$), where $1\leq \nu_1,\nu_2 \leq \frac{p-1}{2}$, having one vertex of $W$, say $b^{\nu_1}_i$ and remaining sets have more than one or do not have any vertices of $W$, then we get the similar contradiction as mentioned in case \ref{item 9}
   \item If all sets $B^{\nu_1}$ and $C^{\nu_2}$ (Where $1\leq \nu_1,\nu_2 \leq \frac{p-1}{2})$ either contain more than one vertex of $W$ or do not contain any vertex of $W$, then we encounter the similar contradiction as mentioned in the case \ref{item 9}.
      \end{enumerate} 
 \end{itemize}
\item  If $q-2$ vertices of $W$ are chosen from $A$, $B$, and $Q$ (or from $A$, $C$, and $Q$, or from $B$, $C$, and $Q$, or from $A$, $B$, $C$, and $Q$ ), in any combination, then we get the same contradiction as mentioned in the case \ref{item1}.
\end{enumerate}
 It has been observed that when $p+1< q < 2p-2$, selecting a set (denoted as $W$) of $q-2$ vertices from $A$, $B$, $C$, and $Q$ in any combination, results in at least two vertices have the same metric coordinates with respect to $W$. This indicates that the MD of  $BS(\Gamma(R))$ is greater than $q-2$ when $p+1< q < 2p-2$.
 \end{proof}

\end{document}
